% =====================================================================
% 10_risksharing_specialness.tex — efficient risk sharing / LOOP,
% the specialness of safe dollar bonds (main paper II.D / II.C),
% Foreign-only demand (App. B.4), distinct capital + sticky prices
% (App. B.5) (WP-E)
% =====================================================================
\section{Efficient risk sharing, the law of one price, and the specialness of safe dollar bonds}
\label{sec:risksharing-specialness}

The previous sections built the analytical machine: the period-$1$ engine
(Section~\ref{sec:engine}), Propositions~\ref{prop:prices}--\ref{prop:wealth} on
prices/returns/wealth, the second-order portfolio engine
(Section~\ref{sec:portfolio-engine}), and Propositions~\ref{prop:portfolios}--%
\ref{prop:riskpremium} on portfolios and the risk premium. This section collects
the three remaining analytical statements of KL's Section~II and the two
robustness exercises of their online Appendix~B. First, we make explicit the
\emph{efficient-risk-sharing / law-of-one-price} content of the model: the
result, stated cleanly in II.C's ``\emph{Efficient Risk Sharing and the Law of
One Price}'' paragraph (their p.~1667), that the dollar appreciation in bad times
is a \emph{necessary consequence} of efficient risk sharing plus the law of one
price, manifested as a tight link between the real exchange rate and relative
output. Second, we transcribe and explain KL's three ``specialness of safe dollar
bonds'' arguments (their II.D, p.~1666). Third, we reconstruct the
Foreign-only-demand robustness (their App.~B.4) and its multiplier result
$\mu_t=\cy_t$. Fourth, we summarize compactly the distinct-capital-stocks-and-%
sticky-prices environment (their App.~B.5), stating the analogs of
Propositions~1--5 and flagging that KL exclude their proofs.

\subsection{The efficient-risk-sharing / law-of-one-price relation}
\label{sec:loop}

\paragraph{From the kernel relation to a relation in quantities.} The
local-completeness condition~\eqref{eq:risk-sharing} of
Section~\ref{sec:local-completeness} is the Backus--Smith / law-of-one-price
relation in \emph{kernel} form,
\begin{equation}
\dev{\pk}_{0,1} - \rra\,\tel^z\dev{\epsilon}^z_1
= \fr{\dev{\pk}}_{0,1} - \rraF\,\tel^z\dev{\epsilon}^z_1 + \D\dev{\rer}_1.
\tag{\ref{eq:risk-sharing} revisited}
\end{equation}
KL observe (their p.~1667) that this relation has a sharp counterpart in
\emph{quantities} once one specializes to the unitary-risk-aversion benchmark
$\rra=\rraF=1$, which --- recalling the maintained unit IES of the simplified
environment --- makes preferences time-separable. In that benchmark the kernel is
the pure consumption-growth term, $\dev{\pk}_{0,1}=-\dev{\rs c}_1$ and
$\fr{\dev{\pk}}_{0,1}=-\fr{\dev{\rs c}}_1$ (the continuation-value term in
\eqref{eq:m01} carries the factor $1-\rra=0$), and the productivity-loading terms
$\rra\tel^z\dev{\epsilon}^z_1$, $\rraF\tel^z\dev{\epsilon}^z_1$ are equal. The
risk-sharing relation~\eqref{eq:risk-sharing} then reduces to the
\emph{Backus--Smith relation in levels} (their displayed equation, p.~1667,
``assuming no shocks prior to period $t$''):
\begin{equation}
\fr{\dev{\rs c}}_1 - \dev{\rs c}_1 = \dev{\rer}_1, \label{eq:bs-levels}
\end{equation}
i.e.\ Home consumption rises relative to Foreign exactly when the Home real
exchange rate depreciates ($\dev{\rer}_1<0$ raises $\dev{\rs c}_1$ relative to
$\fr{\dev{\rs c}}_1$). This is the standard complete-markets prediction that
relative consumption tracks the real exchange rate.
\gapflag{KL display the benchmark relation~\eqref{eq:bs-levels} as
``$\hat c^{*}_t-\hat c_t=\hat q_t$'' for the $\rra=\rraF=1$ case (their p.~1667).
We supply the one-line provenance: it is the local-completeness
condition~\eqref{eq:risk-sharing} evaluated at $\rra=\rraF=1$, where the
continuation-value term of the kernel~\eqref{eq:m01} drops (its coefficient
$1-\rra$ vanishes) and the symmetric productivity loadings cancel, leaving
$-\dev{\rs c}_1=-\fr{\dev{\rs c}}_1+\dev{\rer}_1$, i.e.\
$\fr{\dev{\rs c}}_1-\dev{\rs c}_1=\dev{\rer}_1$. KL's sign convention writes
$\hat q$ as a Home appreciation (Section~\ref{sec:aux-vars},
eq.~\eqref{eq:rer-def}); their displayed line $\hat c^{*}_t-\hat c_t=\hat q_t$ is
exactly~\eqref{eq:bs-levels} in our notation.}

\paragraph{Linking the real exchange rate to relative output.} KL then combine
\eqref{eq:bs-levels} with the engine's relative-goods-market and
intratemporal-allocation relations to obtain a relation between the real exchange
rate and \emph{relative output}. The relevant engine inputs are the
real-exchange-rate--terms-of-trade relation~\eqref{eq:q1-s1},
$\dev{\rer}_1=\hbias\frac{1+\zs}{\zs}\dev{\tot}_1$, and the cross-country
goods-market / intratemporal condition. KL's displayed combination (their
p.~1667, with an \emph{arbitrary} degree of home bias $\hbias$ retained ``to
emphasize the generality of the result'') reads
\begin{equation}
\hbias\,\frac{1+\zs}{\zs}\bigl(\fr{\dev{\rs c}}_1-\dev{\rs c}_1\bigr)
+\Bigl(1+\hbias\,\tfrac{1+\zs}{\zs}\Bigr)
\Bigl(1-\hbias\,\tfrac{1+\zs}{\zs}\Bigr)
\frac{1}{\hbias\,\frac{1+\zs}{\zs}}\,\tel\,\dev{\rer}_1
= \fr{\out}_1 - \out_1, \label{eq:loop-rel-raw}
\end{equation}
where $\out_1,\fr{\out}_1$ are Home and Foreign log output deviations. Using the
Backus--Smith relation~\eqref{eq:bs-levels} to replace
$\fr{\dev{\rs c}}_1-\dev{\rs c}_1=\dev{\rer}_1$ in the first term, the two
left-hand terms collapse into a single multiple of $\dev{\rer}_1$, and solving
for $\dev{\rer}_1$ yields KL's displayed law-of-one-price relation (their
p.~1667):
\begin{equation}
\boxed{\;
\dev{\rer}_1 = \frac{1}
{\hbias\,\frac{1+\zs}{\zs}
+\Bigl(1+\hbias\,\frac{1+\zs}{\zs}\Bigr)\Bigl(1-\hbias\,\frac{1+\zs}{\zs}\Bigr)
\dfrac{1}{\hbias\,\frac{1+\zs}{\zs}}\,\tel}\,
\bigl(\fr{\out}_1 - \out_1\bigr).\;} \label{eq:loop-solved}
\end{equation}
\gapflag{Equations~\eqref{eq:loop-rel-raw}--\eqref{eq:loop-solved} are KL's two
displayed lines on p.~1667 (the relation in
$\varsigma\frac{1+\zeta^{*}}{\zeta^{*}}(\hat c^{*}_t-\hat c_t)$ and $\sigma\hat q_t$
equal to $\hat y^{*}_t-\hat y_t$, and the solved expression for $\hat q_t$). We
transcribe them in the document's macros (KL's
$\varsigma\frac{1+\zeta^{*}}{\zeta^{*}}$ is our $\hbias\frac{1+\zs}{\zs}$, the
recurring relative-price loading of~\eqref{eq:q1-s1}; their $\sigma$ is our
$\tel$). The first line is the log-linearized statement of ``intratemporal
optimality between Home- and Foreign-produced goods in each country, together with
goods market clearing''; it is the same content as the engine's cross-country
goods-market relation~\eqref{eq:rc-rel}, re-expressed in output rather than
terms-of-trade units and with the relative-consumption term written using
$\hbias\frac{1+\zs}{\zs}(\fr{\dev{\rs c}}_1-\dev{\rs c}_1)$. We did not re-derive
the output-form of~\eqref{eq:rc-rel} from primitives; we transcribe KL's
displayed relation and verify that substituting~\eqref{eq:bs-levels} collapses
it to~\eqref{eq:loop-solved}.}

\paragraph{Reading the result.} The denominator of~\eqref{eq:loop-solved} is
strictly positive whenever there is consumption home bias, $\hbias>0$ (so that
$0<\hbias\frac{1+\zs}{\zs}<1$ and the middle factor
$(1-\hbias\frac{1+\zs}{\zs})>0$). Hence~\eqref{eq:loop-solved} delivers KL's
central qualitative statement (their p.~1667):
\begin{quote}
\emph{Given consumption home bias $\hbias>0$, any shock that appreciates the Home
real exchange rate ($\dev{\rer}_1>0$) must be reflected in a relative fall in
Home output ($\out_1<\fr{\out}_1$).}
\end{quote}
KL identify the two distinct channels behind this link, both visible
in~\eqref{eq:loop-rel-raw}:
\begin{enumerate}[leftmargin=1.8em,itemsep=3pt]
  \item \textbf{Efficient risk sharing.} The Backus--Smith
  relation~\eqref{eq:bs-levels} forces relative Home \emph{consumption} to fall
  when the Home real exchange rate appreciates. This is the
  $\hbias\frac{1+\zs}{\zs}(\fr{\dev{\rs c}}_1-\dev{\rs c}_1)$ term: efficient
  risk sharing calls for Home consumption to fall relative to Foreign exactly in
  states of Home appreciation.
  \item \textbf{The law of one price (expenditure switching).} Under
  producer-currency pricing (the law of one price), a Home appreciation raises the
  relative price of Home goods --- the associated improvement in Home's terms of
  trade triggers a global expenditure switch \emph{away} from Home-produced goods,
  depressing Home output. This is the $\tel\,\dev{\rer}_1$ term, with the trade
  elasticity $\tel$ governing its strength.
\end{enumerate}
Both channels push Home output down when the dollar appreciates. KL note (their
p.~1667) that relaxing the law of one price --- e.g.\ under sticky local-currency
pricing --- can sever the \emph{second} channel by decoupling the consumer-facing
relative price from the real exchange rate, but the \emph{first} (efficient
risk-sharing) channel survives as long as risk sharing is efficient. This is why
the result is ``a natural consequence of efficient risk sharing and the law of
one price'' rather than of either alone.

\paragraph{Relation to the literature.} KL place this result against three
alternatives (their p.~1667). \emph{Sauzet (2023)} shows that a fall in relative
US output in crises can resolve the reserve-currency paradox following exactly
this logic; KL's contribution is to \emph{identify the shock} --- the flight to
safe dollar assets --- that produces this output pattern endogenously.
\emph{Maggiori (2017)} resolves the paradox instead through an expenditure shift
toward US goods in crises, which operates through the goods-market term
of~\eqref{eq:loop-rel-raw}, allowing the dollar to appreciate even with relative
output unchanged. \emph{Dahlquist et al.\ (2023)} generate a dollar appreciation
from a fall in US output as in Sauzet but augment with deep habits, so the output
decline also shifts relative demand toward US goods. KL's safety-shock mechanism
delivers the output channel without an exogenous taste shift.

\paragraph{The unitary-risk-aversion benchmark and constant wealth share.} The
benchmark $\rra=\rraF=1$ used to derive~\eqref{eq:bs-levels} is the same benchmark
referenced in Propositions~\ref{prop:wealth}--\ref{prop:portfolios}. KL note
(their footnote~58, App.~B.5) that
\begin{quote}
\emph{consistent with eq.~(99) in the baseline simplified environment, unitary
risk aversions $\rra=\rraF=1$ imply that equilibrium financial portfolios hedge
only labor-income risk,}
\end{quote}
and, correspondingly, that efficient risk sharing keeps Home's \emph{overall}
wealth share (inclusive of labor income) constant in response to all shocks. The
analytical content is visible in the international-portfolio
equation~\eqref{eq:eq99}: every right-hand-side term carries either
$(\rraF-\rra)$ or $\bigl(\frac{1}{\zs}(\rraF-1)+(\rra-1)\bigr)$, both of which
vanish at $\rra=\rraF=1$, so the risk-sharing requirement on $\Eone\dev{\wsh}_2$
collapses to the labor-income-hedging terms alone. This is the benchmark in which
``Home's overall wealth share is constant,'' against which Proposition~4's
comparative statics in $\rraF/\rra$ are taken.
\gapflag{The connection between the $\rra=\rraF=1$ benchmark and the constant
overall wealth share is KL's (their footnote~58 and the discussion at p.~1668:
``in the benchmark with identical risk tolerance across countries, efficient risk
sharing calls for the Home wealth share to increase upon a flight to safety
\emph{provided risk aversion exceeds unity}''). We reconstruct the analytical
reason --- that the $(\rraF-\rra)$ and $\bigl(\frac{1}{\zs}(\rraF-1)+(\rra-1)\bigr)$
loadings of~\eqref{eq:eq99} both vanish at $\rra=\rraF=1$, leaving only the
labor-income-hedging terms, so the overall wealth share (financial plus human) is
constant. Note the distinction KL draw: at $\rra=\rraF=1$ the \emph{overall}
wealth share is constant, whereas for $\rra=\rraF>1$ the overall wealth share
\emph{rises} on a flight to safety (the real-exchange-rate-hedging motive of
relative-supply shocks); the $\rra=\rraF=1$ knife-edge is the one in which the
financial portfolio hedges labor income exactly. We flag that we supply this
reading from the structure of~\eqref{eq:eq99} rather than re-deriving the overall
wealth share from a separate human-wealth accounting.}

\subsection{The specialness of safe dollar bonds}
\label{sec:specialness}

KL devote their Section~II.D (p.~1666) to three dimensions in which the demand
for safe dollar bonds is \emph{special} relative to the demand for other assets.
These are stated as arguments / extensions rather than as formal propositions; we
transcribe each with the available analytical content and match KL's level of
detail (expository, with the formal verification deferred to the cited
extensions).

\subsubsection{Zero net supply}
\label{sec:zero-net-supply}

Safe dollar bonds are in \emph{zero net supply}: the outstanding stock is the debt
issued by the Home government (their fiscal rule~\eqref{eq:fiscal-supply}), held
in positive amounts by private agents, so that private and public positions sum to
zero. KL's argument (their p.~1666) is that this is precisely what makes a flight
to safety \emph{contractionary}:
\begin{quote}
\emph{The demand for safe dollar bonds triggers a Keynesian recession because
these assets are in zero net supply.}
\end{quote}
The contrast is with a counterfactual in which the heightened demand were instead
for \emph{capital} (equity). Formally, suppose capital also entered the utility
function and its nonpecuniary value rose on the margin. Then --- because capital
is in \emph{positive} net supply and is a produced, priced asset --- the increase
in its nonpecuniary value would raise the price of capital $\qk_t$, raise
household wealth (existing capital holders are richer), and thereby \emph{raise}
consumption demand and aggregate output. A demand shift toward capital is
\emph{expansionary}. KL note further that relaxing the assumption of a fixed
global capital stock would reinforce this: the demand for capital would also
stimulate output through increased \emph{investment}.

The economic content is the bond-vs-equity distinction. A safety shock that raises
demand for an asset in \emph{zero} net supply cannot be accommodated by a quantity
expansion (there is no net stock to bid up in aggregate); the only margin is the
price of \emph{other} assets and, through the convenience-yield
wedge~\eqref{eq:euler-r} and nominal rigidity, the real interest rate and
employment. The flight to the safe bond therefore manifests as the inability of
the real rate to clear the goods market at full employment --- a Keynesian
recession --- rather than as a wealth-and-investment expansion. KL frame this as
underscoring the importance of distinguishing convenience yields in \emph{bond}
versus \emph{equity} markets, as in Koijen and Yogo (2020), for understanding
their macroeconomic effects.
\gapflag{The zero-net-supply argument is stated by KL as prose (their p.~1666);
they do not write out the capital-in-utility counterfactual formally. We supply
the mechanism --- that a positive-net-supply, priced asset (capital) absorbs a
demand shift through its price $\qk_t$ and a wealth/investment expansion, whereas
a zero-net-supply asset (the safe bond) can only transmit the shift through the
convenience-yield wedge and nominal rigidity, producing the Keynesian recession of
Proposition~\ref{prop:prices}. This is consistent with, but not a re-derivation
of, KL's argument; the formal capital-in-utility extension is not carried out by
KL and we do not carry it out here. We flag that the ``stimulating'' sign of a
capital-demand shift is asserted, following KL's appeal to Koijen and Yogo (2020).}

\subsubsection{Country size}
\label{sec:country-size}

The second dimension is the \emph{large size} of the US economy. KL argue (their
p.~1666) that the demand for safe dollar bonds is special relative to the demand
for the bonds of \emph{small} issuers. The thought experiment is to augment the
two-country model with an \emph{infinitesimal} third country --- for concreteness,
Switzerland --- that is a negligible part of the global economy. An increase in
the portfolio demand for Swiss-franc bonds would:
\begin{itemize}[leftmargin=1.8em,itemsep=2pt]
  \item generate an \emph{appreciation} of the franc;
  \item cause a \emph{decline} in Swiss production (the same
  efficient-risk-sharing-plus-LOOP logic of Subsection~\ref{sec:loop}: an
  appreciation of the small country's real exchange rate is reflected in a
  relative fall in its output); but
  \item have \emph{negligible spillovers} on global demand and production, because
  the country's share of the world economy is infinitesimal.
\end{itemize}
By contrast, the US is large enough that a flight to safe dollar bonds moves
\emph{global} aggregates --- world employment and output --- not merely US relative
output. KL connect this to the message of Hassan (2013) that the United States'
relative size in the global economy may be an important contributor to the
dollar's negative beta: it is precisely because the dollar issuer is large that
the dollar's appreciation in bad times has aggregate, rather than purely relative,
consequences and hence commands a sizeable (negative-beta) risk premium.
\gapflag{The country-size argument is stated by KL as a thought experiment (their
p.~1666); the infinitesimal-third-country model is described but not formally
written out. We state it exactly as KL do, noting that the ``appreciation $+$
relative-output decline $+$ negligible global spillover'' conclusion for the small
country follows from applying the law-of-one-price / efficient-risk-sharing
relation~\eqref{eq:loop-solved} to a country of vanishing global weight. We do not
formalize the three-country extension; KL do not, and we match their expository
level. The link to Hassan (2013) is KL's.}

\subsubsection{Currency of invoicing}
\label{sec:invoicing}

The third dimension is the widespread use of the dollar in \emph{pricing}. KL
argue (their p.~1666) that dollar invoicing renders the demand for safe dollar
bonds special relative to the currencies of other large countries. The concrete
example: suppose nominal wages \emph{even in Foreign} are denominated and sticky in
dollars (rather than in Foreign currency). Then the global decline in employment
following a positive safety shock is \emph{exacerbated} relative to the baseline
model. The intuition is that dollar deflation following a positive safety shock
(the safety shock is deflationary in the dollar bloc,
Proposition~\ref{prop:prices}) now implies that product wages \emph{in both
countries} are too high --- the sticky dollar wage does not adjust down with the
Foreign price level --- generating a more severe and more uniform global recession.

KL note that similar results obtain under dollar pricing of \emph{exports} (the
dominant-currency-pricing paradigm of Gopinath (2015), Gopinath et al.\ (2020),
and Mukhin (2022)): dollar invoicing of trade among countries not involving the
United States is analogous, in the two-country setup, to an internal Foreign price
(such as the nominal wage) being denominated in dollars. The general message is
that the dollar's role in goods markets (invoicing) and its role in financial
markets (safe-asset demand) interact: dollar invoicing amplifies the employment
consequences of dollar safe-asset demand, suggesting ``potentially rich
interactions between the dollar's role in financial and goods markets.''
\gapflag{The currency-of-invoicing argument is stated by KL as prose (their
p.~1666); the dollar-wage and dollar-export-pricing extensions are described but
not formally derived. We transcribe KL's argument and its mapping (dollar
invoicing of trade $\approx$ an internal Foreign price denominated in dollars, in
the two-country reduction). We do not formalize either extension; KL state them as
suggestive, citing Gopinath (2015), Gopinath et al.\ (2020), and Mukhin (2022) for
the dominant-currency-pricing channel. We match their expository level.}

\subsection{Robustness: Foreign-only demand for safe dollar bonds (App.~B.4)}
\label{sec:foreign-only}

KL's online Appendix~B.4 shows that the analysis is robust to the realistic case
in which \emph{only Foreign} households have a special (nonpecuniary) demand for
safe dollar bonds, with a non-negativity constraint preventing private agents from
\emph{creating} safe dollar assets (only the Home government can). We reconstruct
the FOC algebra, which is the load-bearing part of B.4.

\paragraph{The augmented household FOCs.} KL augment the model with a
non-negativity constraint on each household's position in the safe dollar bond,
reflecting that households cannot manufacture safe assets. This is irrelevant when
the demand shock is \emph{global} (as in the baseline), since then all agents hold
a positive share of the government's outstanding safe debt; it binds when the
demand is \emph{Foreign-only}, because Home households then have no nonpecuniary
motive to hold the safe bond and would, absent the constraint, want to short it.
The Home representative agent's FOCs for safe dollar bonds, other dollar bonds,
Foreign bonds, and capital become (their app.~p.~79)
\begin{align}
1 - \mu_t &= \Et\,\pk_{t,t+1}(1+\inom_t)\frac{P_t}{P_{t+1}}, \label{eq:b4-home-safe}\\
1 &= \Et\,\pk_{t,t+1}(1+\ioth_t)\frac{P_t}{P_{t+1}}, \label{eq:b4-home-other}\\
1 &= \Et\,\pk_{t,t+1}\frac{\rer_t}{\rer_{t+1}}(1+\fr{r}_{t+1}), \label{eq:b4-home-fbond}\\
1 &= \Et\,\pk_{t,t+1}(1+\rk_{t+1}), \label{eq:b4-home-cap}
\end{align}
where $\mu_t\ge0$ is the (scaled) multiplier on the non-negativity constraint on
the Home agent's safe-dollar-bond position, and the last three FOCs are exactly as
in the baseline model. The Foreign representative agent --- who \emph{does} receive
the nonpecuniary safe-bond utility, with the same functional form $\bondu^{*}_t$ as
in the baseline --- has FOCs (their app.~p.~79)
\begin{align}
1 - \fr c_t\,\frac{\bondu^{*}_t{}'\!\bigl(\fr B_{Ht,s}/(E_t^{-1}\fr P_t)\bigr)}
{\bondu^{*}_t\!\bigl(\fr B_{Ht,s}/(E_t^{-1}\fr P_t)\bigr)}
&= \Et\,\fr{\pk}_{t,t+1}\frac{\rer_{t+1}}{\rer_t}(1+\inom_t)\frac{P_t}{P_{t+1}}, \label{eq:b4-foreign-safe}\\
1 &= \Et\,\fr{\pk}_{t,t+1}\frac{\rer_{t+1}}{\rer_t}(1+\ioth_t)\frac{P_t}{P_{t+1}}, \label{eq:b4-foreign-other}\\
1 &= \Et\,\fr{\pk}_{t,t+1}(1+\fr{r}_{t+1}), \label{eq:b4-foreign-fbond}\\
1 &= \Et\,\fr{\pk}_{t,t+1}\frac{\rer_{t+1}}{\rer_t}(1+\rk_{t+1}), \label{eq:b4-foreign-cap}
\end{align}
as in the baseline model. Given the same functional form for $\bondu^{*}_t$ as in
the baseline (eq.~\eqref{eq:Omega-form}, starred), the Foreign marginal
nonpecuniary value evaluates to
\begin{equation}
\fr c_t\,\frac{\bondu^{*}_t{}'\!\bigl(\fr B_{Ht,s}/(E_t^{-1}\fr P_t)\bigr)}
{\bondu^{*}_t\!\bigl(\fr B_{Ht,s}/(E_t^{-1}\fr P_t)\bigr)}
= \omega^{d*}_t - \frac{\fr B_{Ht,s}}{E_t^{-1}\fr P_t\,\fr c_t}, \label{eq:b4-foreign-wedge}
\end{equation}
where $\omega^{d*}_t$ is the Foreign latent demand shock for safe dollar bonds. KL
assume for expositional simplicity that the non-negativity constraint for
\emph{Foreign} agents does not bind, which is the case when $\omega^{d*}_t$ is
sufficiently high.

\paragraph{The multiplier equals the convenience yield.} The key step. Each
agent's pair of FOCs for safe and other dollar bonds imply, by dividing the
safe-bond FOC by the other-bond FOC. For the Home agent,
\eqref{eq:b4-home-safe}/\eqref{eq:b4-home-other} gives
\begin{equation}
\frac{1+\inom_t}{1-\mu_t} = 1+\ioth_t, \label{eq:b4-home-ratio}
\end{equation}
and for the Foreign agent,
\eqref{eq:b4-foreign-safe}/\eqref{eq:b4-foreign-other} gives, using
\eqref{eq:b4-foreign-wedge},
\begin{equation}
1+\ioth_t = \frac{1+\inom_t}{1-\Bigl(\omega^{d*}_t-\dfrac{\fr B_{Ht,s}}{E_t^{-1}\fr P_t\,\fr c_t}\Bigr)}. \label{eq:b4-foreign-ratio}
\end{equation}
Both expressions equal the common gross other-dollar-bond return $1+\ioth_t$, so
their denominators must coincide:
\begin{equation}
\frac{1+\inom_t}{1-\mu_t}
= \frac{1+\inom_t}{1-\Bigl(\omega^{d*}_t-\dfrac{\fr B_{Ht,s}}{E_t^{-1}\fr P_t\,\fr c_t}\Bigr)}
\;\Longrightarrow\;
\boxed{\;\mu_t = \omega^{d*}_t - \frac{\fr B_{Ht,s}}{E_t^{-1}\fr P_t\,\fr c_t} \equiv \cy_t.\;} \label{eq:b4-mu-omega}
\end{equation}
That is, the Lagrange multiplier on the non-negativity constraint for safe dollar
bonds faced by Home agents must equal the convenience yield $\cy_t$ perceived by
Foreign agents on those bonds. Intuitively, the shadow value of being unable to
short the safe bond (for Home) is exactly the convenience premium that Foreign
agents attach to it.
\gapflag{Equations~\eqref{eq:b4-home-ratio}--\eqref{eq:b4-mu-omega} reconstruct
KL's App.~B.4 result
$\mu_t=\omega^{d*}_t-\fr B_{Ht,s}/(E_t^{-1}\fr P_t\fr c_t)\equiv\cy_t$ (their
app.~p.~79--80). KL display the FOCs~\eqref{eq:b4-home-safe}--%
\eqref{eq:b4-foreign-cap}, the Foreign wedge~\eqref{eq:b4-foreign-wedge}, the
intermediate identity
$\frac{1+i_t}{1-\mu_t}=1+\ioth_t=\frac{1+i_t}{1-(\omega^{d*}_t-\fr B_{Ht,s}/(E_t^{-1}\fr P_t\fr c_t))}$,
and conclude
$\mu_t=\omega^{d*}_t-\fr B_{Ht,s}/(E_t^{-1}\fr P_t\fr c_t)\equiv\cy_t$. We supply the
one-line division step (each agent's safe-bond FOC over its other-bond FOC
produces the gross convenience-adjusted return, and equating the two expressions
for $1+\ioth_t$ forces the denominators to coincide). The definition
$\cy_t\equiv\omega^{d*}_t-\fr B_{Ht,s}/(E_t^{-1}\fr P_t\fr c_t)$ is the Foreign-only
analog of the baseline reduced-form convenience yield~\eqref{eq:omega-reduced},
now reflecting only the Foreign latent demand $\omega^{d*}_t$ (and the supply of safe
bonds by the Home government). We reproduce KL's algebra faithfully.}

\paragraph{Consequence: the equilibrium system reduces to the baseline with
$b_{Ht,s}=0$.} With the multiplier pinned by~\eqref{eq:b4-mu-omega}, the Home
agent is at the corner --- it holds \emph{zero} safe dollar bonds, $b_{Ht,s}=0$ ---
and the Home safe-bond FOC~\eqref{eq:b4-home-safe} reads
$1-\cy_t=\Et\pk_{t,t+1}(1+\inom_t)P_t/P_{t+1}$, which is exactly the baseline
convenience-adjusted Euler equation~\eqref{eq:euler-safe} with the Home agent
holding none of the safe asset. KL conclude (their app.~p.~80):
\begin{quote}
\emph{The equilibrium system is thus exactly as in the baseline model, except with
$B_{Ht,s}=0$ (zero holdings of safe dollar bonds by Home agents). It follows that
the propagation of Foreign demand shocks for safe dollar bonds is exactly like
global demand shocks for these assets, up to any differences in seigniorage
because in this case only Foreign agents hold them.}
\end{quote}
In particular, when the Home government issues zero net safe debt ($b^g_{H,s}=0$)
there is no seigniorage, and the propagation of Foreign demand shocks is
\emph{identical} to that of global demand shocks (their footnote~55). A Foreign
demand shift propagates exactly like a baseline safety shock: at unchanged policy
rates, Foreigners seek to borrow in non-safe dollar bonds to hold safe dollar
bonds (capturing the convenience benefit without changing their currency
exposure), which raises the other-dollar-bond rate $\ioth_t$ in the absence of a
decline in the policy rate $\inom_t$, making saving more attractive for Home
households and inducing the same effects on quantities and relative prices as in
the baseline model.

\paragraph{The international risk-sharing conditions are unchanged.} KL emphasize
(their app.~p.~80) that despite only Foreigners having a special demand for safe
dollar bonds, the \emph{international risk-sharing conditions are unchanged}. The
three first-order risk-sharing relations remain
\begin{align}
\Et\,\pk_{t,t+1}\frac{1+\rr_{t+1}}{1-\cy_t}
&= \Et\,\fr{\pk}_{t,t+1}\frac{\rer_{t+1}}{\rer_t}\frac{1+\rr_{t+1}}{1-\cy_t}, \label{eq:b4-rs-safe}\\
\Et\,\pk_{t,t+1}\frac{\rer_t}{\rer_{t+1}}(1+\fr{r}_{t+1})
&= \Et\,\fr{\pk}_{t,t+1}(1+\fr{r}_{t+1}), \label{eq:b4-rs-fbond}\\
\Et\,\pk_{t,t+1}(1+\rk_{t+1})
&= \Et\,\fr{\pk}_{t,t+1}\frac{\rer_{t+1}}{\rer_t}(1+\rk_{t+1}), \label{eq:b4-rs-cap}
\end{align}
i.e.\ the same Home/Foreign kernel-equality conditions (up to the
real-exchange-rate conversion) that underlie the local-completeness
relation~\eqref{eq:risk-sharing}. With one more asset than shocks, the simplified
environment remains \emph{locally complete} as in the prior subsections, so
Propositions~\ref{prop:prices}--\ref{prop:riskpremium} carry over (with
$b_{Ht,s}=0$ on the Home side and the convenience yield now reflecting Foreign
demand only).
\gapflag{Equations~\eqref{eq:b4-rs-safe}--\eqref{eq:b4-rs-cap} transcribe KL's
displayed ``international risk sharing conditions'' (their app.~p.~80), which they
note are unchanged from the baseline. We reproduce them and the conclusion that
local completeness (one more asset than shocks) is preserved. The content
load-bearing for the big\_cy reading is exactly KL's statement: Foreign-only
safe-dollar demand propagates like a global safety shock (up to seigniorage), and
the risk-premium/portfolio results are unchanged --- so the baseline analytical
results are robust to the more realistic Foreign-only demand structure. We did not
re-derive the propagation; KL assert it follows from the system reducing to the
baseline with $B_{Ht,s}=0$, and we transcribe that reduction.}

\subsection{Distinct capital stocks and sticky prices (App.~B.5)}
\label{sec:b5}

KL's online Appendix~B.5 studies an \emph{alternative} simplified environment with
(i) \emph{distinct} national capital stocks --- Home agents choose separate
positions $k_{Ht}$ (capital used at Home) and $k_{Ft}$ (capital used at Foreign),
which offer distinct payoffs and trade at distinct prices --- and (ii) \emph{sticky
prices} (one-period-in-advance \emph{price} setting by monopolistically
competitive retailers) in place of the baseline's sticky \emph{wages}. Because
there is now trade in one more asset (two distinct capital claims rather than a
single global capital stock), one more shock is needed to pin down portfolios; KL
add a fully transitory Foreign productivity shock $z_{Ft}$ (with $\rho^{F}=0$),
analogous to the safety shock. KL provide analogs of Propositions~1--5 in this
environment but \emph{exclude the proofs} (``available on request''). We summarize
the results compactly and flag the omitted proofs.

\gapflag{This entire subsection summarizes KL's App.~B.5, which KL state
\emph{without proof} (their app.~p.~83: ``We now provide analogs of
Propositions~1--5 in this environment. We exclude the proofs for brevity but they
are available on request.''). We therefore state the analog results as KL display
them and do \emph{not} reconstruct their proofs --- a full re-derivation is out of
scope for this faithful reference and is not available in the published material.
Where a result differs qualitatively from the baseline (most importantly the
ambiguous sign of the Home capital return under sticky prices), we note the
difference and KL's stated intuition. The big\_cy extension worker should treat
B.5 as a stated-but-unproved alternative environment.}

\paragraph{Environment (B.5.1).} Relative to the baseline (Section~II): Home
households hold distinct claims $k_{Ht},k_{Ft}$ to Home- and Foreign-used capital,
with fixed national stocks $\bar k,\bar k^{*}$. Consumption of each country's
goods is a CES aggregate across varieties with elasticity $\lel$. Each country has
monopolistically competitive intermediate-good producers and final-good retailers;
retailers set prices either flexibly or one period in advance (the sticky-price
analog of the baseline's sticky wages). Households own a share of all retailers
from a given country in proportion to their ownership of that country's capital, so
aggregate profits accrue to capital owners. Capital-market clearing is now
$k_{Ht}+\zs\fr k_{Ht}=\bar k$ and $k_{Ft}+\zs\fr k_{Ft}=\bar k^{*}$. KL retain the
maintained simplifications of the baseline (e.g.\ no disaster risk or depreciation
in the analytical treatment) but \emph{do not} impose the infinite Frisch
elasticity. Because markets remain \emph{locally complete}, the results are
independent of which capital owner's stochastic discount factor is used to price
firms' dynamic decisions.

\paragraph{Proposition~1 analog (prices and quantities).} With
$\Et\dev{\wsh}_{t+1}=0$ on impact of all shocks (the symmetric-portfolio benchmark
of Propositions~\ref{prop:prices}--\ref{prop:wealth}), the transmission of a
positive safety shock is \emph{largely unchanged} from
Proposition~\ref{prop:prices}. Under flexible prices: the Home CPI declines
($\D\dev P_t=-\frac{1}{\taylor}\dev{\cy}_t$), the Home real interest rate declines
($\Et\rr_{t+1}=-\dev{\cy}_t$), and the Home real exchange rate and employment in
each country are unchanged ($\dev{\rer}_t=\dev{\ell}_t=\fr{\dev{\ell}}_t=0$). Under
prices set one period in advance: the Home CPI is unchanged ($\D\dev P_t=0$), the
Home real interest rate is unchanged ($\Et\rr_{t+1}=0$), the Home real exchange
rate \emph{appreciates} ($\dev{\rer}_t=\dev{\cy}_t$), and global employment falls,
disproportionately so in Home
($\frac{1}{1+\zs}\dev{\ell}_t+\frac{\zs}{1+\zs}\fr{\dev{\ell}}_t
=-\frac{1}{1-\cshare}\dev{\cy}_t$ and
$\dev{\ell}_t-\fr{\dev{\ell}}_t=-\frac{1}{1-\cshare}\dev{\cy}_t$). KL conclude the
transmission of safety shocks to prices and quantities is largely unchanged from
Proposition~\ref{prop:prices}.

\paragraph{Proposition~2 analog (returns), with the ambiguous-sign result.}
Defining the real returns on the two distinct capital stocks
$1+\rk_t=(\divr_t+\qk_t)/\qk_{t-1}\cdot P_{t-1}/P_t$ and the Foreign analog
$1+r^{k*}_t$, the analog of Proposition~\ref{prop:returns} on impact of a
positive safety shock is: under flexible prices, the dollar-bond real return rises
($\rr_t=\frac{1}{\taylor}\dev{\cy}_t$) while the Foreign-bond, Home-capital, and
Foreign-capital returns are unaffected (each zero net of the real-exchange-rate
change); under prices set one period in advance, the dollar-bond return is
unchanged ($\rr_t=0$), the Foreign-bond return falls
($\fr{r}_t-\D\dev{\rer}_t=-\dev{\cy}_t$), the Foreign-capital return falls
($r^{k*}_t-\dev{\rer}_t=-\dev{\cy}_t$), and --- the key new feature --- the
\emph{Home}-capital return has an \emph{ambiguous} sign:
\begin{equation}
\rk_t = \left[\frac{(1-\disc)(1-\lwedge)\bigl(\tfrac{1}{\frisch}+1\bigr)}{1-(1-\cshare)(1-\lwedge)}-1\right]\dev{\cy}_t. \label{eq:b5-rk-ambiguous}
\end{equation}
KL explain (their app.~p.~84) that the ambiguous response of the Home capital
return owes to \emph{competing effects} of a safety shock under sticky prices: the
shock \emph{reduces Home output} (pushing the capital return down), but it
\emph{raises the Home markup} $P_{Ht}/P^{i}_t$ (pushing the capital return up,
since with sticky output prices and a falling marginal cost the price--cost margin
widens). What is \emph{definitive}, however, is that the real return on Home
capital \emph{exceeds} that on Foreign capital:
\begin{equation}
\rk_t - r^{k*}_t + \dev{\rer}_t \;\propto\; \dev{\cy}_t \;>\;0 \quad\text{(on a positive safety shock).} \label{eq:b5-rk-gap}
\end{equation}
This relative-capital-return result plays the role in B.5 that the
positive-excess-return results play in the baseline: Home capital outperforms
Foreign capital on a flight to safety. KL note that the markup-rise channel --- a
positive safety shock raising relative Home equity returns through higher markups
under sticky prices --- ``can help to account for the rise in relative US equity
returns in crises documented in Dahlquist et al.\ (2023).''

\paragraph{Proposition~3 analog (wealth share / portfolios at $\rra=\rraF=1$).}
With prices set one period in advance and the unitary-risk-aversion benchmark
$\rra=\rraF=1$, the equilibrium portfolios are
\begin{equation}
\frac{\qk k_H}{\sav}=1,\quad \frac{\qk k_F}{\sav}=0,\quad
\frac{b_F}{\sav}=-\disc\,\frac{\zs}{1+\zs}\,\frac{b^g_{H,s}}{\sav},\quad
\frac{b_H}{\sav}=\disc\,\frac{\zs}{1+\zs}\,\frac{b^g_{H,s}}{\sav}, \label{eq:b5-prop3-port}
\end{equation}
i.e.\ Home fully owns its domestic capital stock and holds none of Foreign capital,
and the dollar-bond/Foreign-bond positions are the seigniorage-hedging pair. These
imply that Home's \emph{financial} wealth share $\dev{\wsh}_t$ \emph{rises} on
impact of a positive safety shock (because, by~\eqref{eq:b5-rk-gap}, Home capital
outperforms Foreign capital), \emph{falls} on a positive Foreign-productivity
shock, and is unchanged on a global-productivity shock; while Home's net foreign
assets $\nfa_t$ and \emph{expected} future wealth share $\Et\dev{\wsh}_{t+1}$ are
unchanged on impact of all shocks. KL emphasize (their footnote~58 and p.~85) that
in this $\rra=\rraF=1$ benchmark, efficient risk sharing calls for Home's
\emph{overall} wealth share (inclusive of labor income) to be \emph{constant} in
response to all shocks: with full domestic capital ownership, financial returns
exactly offset labor-income changes in response to safety and Foreign-productivity
shocks, and wealth does not redistribute across countries in response to
global-productivity shocks (building on Engel and Matsumoto 2009 and Coeurdacier
and Gourinchas 2016). A notable implication: even though Home's overall wealth
share and NFA \emph{fall} on a safety shock (Home is insuring Foreign), its
\emph{financial} wealth share $\dev{\wsh}_t$ can \emph{rise} on impact, because the
Home capital return outperforms Foreign capital.

\paragraph{Proposition~4 analog (portfolios around the benchmark).} Around the
$\rra=\rraF=1$, $b^g_{H,s}<0$ benchmark with a common positive steady-state labor
wedge, the comparative statics mirror Proposition~\ref{prop:portfolios}: Home's
portfolio share in Home capital (Foreign capital, dollar bonds) is unaffected
(unaffected, falls) with $-b^g_{H,s}$; and rises (rises, falls) with $\rraF/\rra$
holding $\rra+\frac{1}{\zs}\rraF$ fixed. That is, efficient risk sharing calls for
Home's overall wealth share to \emph{fall} on a positive safety shock as Home gets
more risk tolerant than Foreign, implemented by Home owning a \emph{levered}
portfolio of Home \emph{and} Foreign capital financed by dollar bonds.

\paragraph{Proposition~5 analog (risk premium).} KL state that the final result
characterizes the currency risk premium around the $\rra=\rraF=1$, $b^g_{H,s}<0$
benchmark and is ``essentially identical to Proposition~\ref{prop:riskpremium}'':
the covariance $\Cov_t(-\dev{\pk}_{t,t+1},\fr{r}_{t+1}-\D\dev{\rer}_{t+1}-\rr_{t+1})$
is proportional to $\rra-\rraF$ when $\sigma^{\cy}=0$ and is rising in
$\sigma^{\cy}$. The reserve-currency-paradox logic of
Proposition~\ref{prop:riskpremium} therefore survives intact in the
distinct-capital, sticky-price environment.

\paragraph{Where B.5 differs from the baseline.} Two qualitative differences are
worth recording for the extension worker. First, under sticky \emph{prices} (B.5)
the Home capital return responds \emph{ambiguously} to a safety
shock~\eqref{eq:b5-rk-ambiguous} --- competing output-decline and markup-rise
effects --- whereas under sticky \emph{wages} (baseline)
Proposition~\ref{prop:returns} gives a clean realized capital-return response. The
markup-rise channel is genuinely new and is KL's structural account of rising
relative US equity returns in crises. Second, KL note (their p.~1669) that with
sticky wages and distinct capital stocks, ``the returns to capital at Home and
Foreign are still equated in response to safety shocks'' --- capital is not
reallocated across borders even if it could be, because any increase in Foreign
output relative to Home induces a Home real-exchange-rate appreciation that fully
offsets the increase in the Foreign return to capital --- whereas the
relative-return result $\rk_t>r^{k*}_t$~\eqref{eq:b5-rk-gap} that drives the
financial-wealth-share dynamics requires the sticky-\emph{price} environment (and
holds away from the limit of complete home bias provided $\tel\neq1$). KL present
these as showing that the baseline's assumptions on capital mobility are ``not
crucial'' for the results.
